\documentclass{article}
\usepackage{tikz}
\usetikzlibrary{shapes, shapes.misc, shapes.symbols, positioning, calc, arrows.meta}

\begin{document}
\begin{figure}[htbp]
\centering
\vspace{0.5cm}
\begin{tikzpicture}[>=stealth, font=\small, thick, scale=0.8, transform shape]

% -------------------------
% Node Formatting Definitions
% -------------------------
\tikzset{
    act/.style={draw, rounded rectangle, minimum height=0.7cm, fill=white, align=center, thick},
    note/.style={draw, shape=note, fill=white, align=left, font=\scriptsize, inner sep=0.2cm},
    sync/.style={draw, fill=black, rectangle, minimum height=0.15cm, inner sep=0},
    init/.style={fill=black, circle, minimum size=0.4cm, inner sep=0},
    finalouter/.style={draw, circle, minimum size=0.5cm, inner sep=0, thick},
    finalinner/.style={fill=black, circle, minimum size=0.25cm, inner sep=0}
}

% -------------------------
% 1. Central Vertical Axis Nodes
% -------------------------
\node[init] (Start) at (0, 0) {};

\node[act] (UI) at (0, -1.5) {User Interface};
\node[act] (WCF) at (0, -3.5) {Call to MVC Controller for user request};
\node[act] (Deliver) at (0, -5.5) {Deliver call to EF Core business logic};

% -------------------------
% 2. Synchronization Bars
% -------------------------
\node[sync, minimum width=14cm] (Fork) at (0, -7.5) {};
\node[sync, minimum width=14cm] (Join) at (0, -12.5) {};

% -------------------------
% 3. Explanatory Note Blocks
% -------------------------
\node[note] (N1) at (4, -2.5) {Generate analytic event};
\draw[thin] (N1.west) -- (0, -2.5);

\node[note] (N2) at (4, -4.5) {Execute controller logic\\for event};
\draw[thin] (N2.west) -- (0, -4.5);

\node[note] (N3) at (4, -6.5) {Call to specific Entity\\Framework functions};
\draw[thin] (N3.west) -- (0, -6.5);

% -------------------------
% 4. Parallel Process Branches (V-Staggered to prevent cluster overlap)
% -------------------------
\node[act] (T1) at (-5.5, -10.5) {Fetch Category\\Stats Data};
\node[act] (T2) at (-3.3, -9) {Fetch Issue\\Status Data};
\node[act] (T3) at (-1.1, -10.5) {Calculate SLA\\Breaches};
\node[act] (T4) at (1.1, -9) {Count Active\\Staff/Workers};
\node[act] (T5) at (3.3, -10.5) {Fetch Recent\\Reports Data};
\node[act] (T6) at (5.5, -9) {Generate PDF\\Report};

% -------------------------
% 5. Reporting Terminal Nodes
% -------------------------
\node[act] (Generate) at (0, -14.5) {Generate Report and View Data Accordingly};

\node[finalouter] (EndO) at (0, -16) {};
\node[finalinner] (EndI) at (0, -16) {};

% -------------------------
% 6. Path Renderings
% -------------------------
% Main Vertical drops
\draw[->] (Start) -- (UI);
\draw[->] (UI) -- (WCF);
\draw[->] (WCF) -- (Deliver);
\draw[->] (Deliver) -- (Fork);

% Fan Out from Fork
\draw[->] (Fork) -- (T1.north);
\draw[->] (Fork) -- (T2.north);
\draw[->] (Fork) -- (T3.north);
\draw[->] (Fork) -- (T4.north);
\draw[->] (Fork) -- (T5.north);
\draw[->] (Fork) -- (T6.north);

% Converge into Join
\draw[->] (T1.south) -- (Join);
\draw[->] (T2.south) -- (Join);
\draw[->] (T3.south) -- (Join);
\draw[->] (T4.south) -- (Join);
\draw[->] (T5.south) -- (Join);
\draw[->] (T6.south) -- (Join);

% Exit Vertical drops
\draw[->] (Join) -- (Generate);
\draw[->] (Generate) -- (EndO);

\end{tikzpicture}
\caption{Activity Diagram illustrating strictly vertical parallel execution reporting flows}
\end{figure}
\end{document}
